\documentclass[12pt]{article}
\usepackage[utf8]{inputenc}
\usepackage[T1]{fontenc}
\usepackage{hyperref}
\title{Tamnim Citak Language Resources for K-12 Curriculum Development}
\author{Compiled by Claude Sonnet 5 (claude-sonnet-5) for Andr\'as Kornai}
\date{}
\begin{document}
\maketitle

\section{Introduction}
Tamnim Citak (ISO 639-3: tml) is spoken in Mappi Regency, South Papua, Indonesia. Search results disagree on its classification: Wikipedia treats it as a variety of Citak, and an older mirror of Wikipedia treats it as a variety of Asmat. Joshua Project lists about 400 Tamnim Citak people and describes the language as a living first language. Very little material exists that is specific to Tamnim Citak.

\section{Primary Educational Resources}
No primer, textbook or school programme for Tamnim Citak turned up in search results.

\section{Linguistic and Reference Materials}
The only reference found is the general Wikipedia article on the Citak language, which lists Tamnim Citak as a dialect: \url{https://en.wikipedia.org/wiki/Citak_language}. No dictionary, grammar or orthography specific to Tamnim Citak was found.

\section{Religious and Community Resources}
Joshua Project has a people-group profile, which lists a New Testament as Scripture: \url{https://joshuaproject.net/people_groups/11391}. ScriptureEarth lists scripture and audio recordings under the code tml (the site also gives the alternative name Asmat Darat): \url{https://www.scriptureearth.org/00eng.php?iso=tml}. I did not open either page, so I cannot confirm which scripture text is in the language itself.

\section{Digital and Audio Resources}
ScriptureEarth is reported to list audio recordings from Global Recordings Network. No other digital resources were found.

\section{Curriculum Development Context}
Sources say the language is a living first language, and it is not known to be taught in schools. A curriculum would have to start from basic documentation: orthography, word lists and recorded texts. It should be developed with the community, and the classification question (Citak or Asmat) should be settled with linguists first.
\end{document}